\section{Security Incident：从危机到可验证改进}

Identity provider 是高价值目标，安全事件会直接挑战客户最核心的信任假设。课堂强调事件带来的痛苦和文化转型；官方 RCA 则提供一个更具体的系统视角：credential 如何暴露、日志查询为何遗漏、session token 如何被利用、何时 containment、怎样通知客户，以及哪些 control 被永久改变。

\subsection{Incident response 五阶段}

一个完整 incident lifecycle 至少包含 detection、containment、investigation、remediation 和 notification。Detection 可能来自内部 telemetry，也可能来自客户或合作伙伴；containment 要禁用 account、revoke session/token 并限制接口；investigation 建立 timeline 和 affected scope；remediation 修复 credential、logging、workflow 与 product control；notification 则向受影响客户提供可执行信息。

\lecturefigure{08-incident-response.png}{安全事件只有关闭 detection、containment、investigation、remediation 与 notification 学习闭环才算结束。}{本地字幕 15:20--19:10；Okta 2023 support-system RCA；概念重绘。}

\begin{knowledgebox}{读图：RCA 要解释“为什么没更早发现”}
Root cause 不只寻找最初错误，还要解释防线为何未阻止、telemetry 为何未显示、investigation query 为何漏掉、session revoke 为何延迟，以及客户信号如何进入响应。官方 2023 RCA 说明不同文件访问路径生成不同 log event，初始查询因此不完整；这个细节展示 observability schema 本身也属于 security control。
\end{knowledgebox}

\begin{importantbox}{Containment 与 remediation 不同}
禁用 service account、revoke session 和阻断 IP 是 containment，目标是停止当前伤害；阻止个人 profile 保存公司 credential、增强文件访问 telemetry、绑定管理员 session、改进 customer notification 才是 remediation，目标是降低复发概率和缩短未来 detection time。两者需要不同 owner 和完成证据。
\end{importantbox}

\begin{warningbox}{透明度不是一次博客发布}
Customer trust 需要及时 scope、受影响对象、recommended action、已知未知项和后续更新。过早给出未经证实的确定结论会反噬信任，过晚披露又让客户失去 containment 时间。Incident communication 应与 forensic confidence 同步，并保存谁在什么时间知道什么的 decision log。
\end{warningbox}

\teachervoice{McKinnon 说安全事件是职业生涯最艰难的时期之一，因为 headline 会把身份公司与“被攻破”直接绑定。值得保留的课堂判断是：无法靠口号恢复信任，只能把痛苦转成 security-first 的资源、行为和可验证控制。}

\subsection{Security-first culture 是 operating system}

Culture 不是墙上的 value list，而是领导者反复询问什么、launch 时愿意延迟什么、预算给谁、晋升奖励什么、坏消息是否被惩罚。若“security first”只要求 security team 加班，而 product roadmap、sales commitment 和 executive review 不改变，组织接收到的真实信号仍是 feature first。

\lecturefigure{09-security-first-culture.png}{Security-first culture 由领导行为、激励、资源和证据共同实现。}{本地字幕 18:00--21:30；概念重绘。}

\begin{knowledgebox}{读图：Culture 可以被审计}
Leader behavior 看高层是否主动问风险和接受 escalation；incentive 看目标、绩效与 launch criteria；resource 看平台改造、security staffing 和 customer support 是否获得预算；evidence 看 control test、incident metric、audit 和 postmortem 是否进入 review。Culture 虽然难以精确量化，却能通过重复行为和 resource allocation 观察。
\end{knowledgebox}

\subsection{本章小结}

Incident response 的目标不是“恢复绿灯”，而是建立可复查的 learning loop。Security-first 只有进入决策、资源和工程证据，才会在下一次压力中生效。

